\scanpage{29}
% transcription
la catégorie dérivée $D(\operatorname{Ab}(X))$ formée des complexes
\deleted{\unsure{à cohomologie}} \deleted{qui sont} tels que $\underline H^i(K^{\unsure{\cdot|'}})=0$ si $i\notin[0,n]$. Les
opérations image inverse et image directe \deleted{sont} sur les
$n$-champs de Picard strict correspondent aux
opérations analogues dans les catégories dérivées (sauf
qu'il faut tronquer, dans le cas de l'image
directe). Les $\underline{\operatorname{Hom}}$ de cat. de Picard \unsure{stricts|strictes}
correspondent aux $\mathrm R\underline{\operatorname{Hom}}$ tronqués \inserted{sur $[0,n]$} dans les
catégories dérivées.

\uline{NB} Il y aurait lieu d'interpréter de façon
géométrique les objets tronqués de $D(\operatorname{Mod}(O))$,
$O$ un Anneau sur $X$, en termes de cat. de
Picard strictes où $O$ opère, et d'interpréter géométriquement les opérations $\mathrm Rf_*$, \unsure{de|et} $\mathrm Lf^*$ et $\mathrm R\underline{\operatorname{Hom}}$.
Interprétation des $\overset{(\mathrm L)}{\otimes}_{\unsure{k|Z}}$ en termes d'accouplements
\inserted{``universels''} de cat. de Picard ?

10. Etude des $n$-Gr champs. \deleted{\unsure{E}} J'ai regardé le
cas $n=1$, $X$ quelconque\inserted{, cf lettre à Deligne}. Mais \unsure{même} si $X$ est le
topos ponctuel, je ne sais quoi dire pour le
cas $n=2$ p.ex. Il faudrait examiner le cas
de la $n$-\inserted{Gr-}catégorie associée à un groupe $\frac12$ simplicial (p.ex. via un groupe topologique), cf \S~2, et voir
la $n$-catégorie des homomorphismes de \deleted{\unsure{groupes}} \inserted{l'objet}
$\frac12$ simplicial unité dans ledit groupe $\frac12$ simplicial.

11. $n$-\deleted{catégorie}\inserted{champ} de Picard \inserted{(pas strict !)} enveloppant d'un $(n)$-$\otimes$ champ $A$
ACU. Ses invariants $\pi_i$ sont les faisceaux $\underline{K}_i(A)$.
% archive sheet number: 28
